\section{Phân tích What-Why}

\subsection{What -- Phân tích dữ liệu}
Dữ liệu trung tâm là bảng phim, sau đó được mở rộng thành các bảng phụ theo quốc gia, thể loại, đạo diễn và diễn viên. Cấu trúc này phù hợp với nhiều nhiệm vụ trực quan hóa khác nhau:
\begin{itemize}[leftmargin=2em]
    \item \textbf{Biến định danh và mô tả:} \texttt{show\_id}, \texttt{title}, \texttt{description}.
    \item \textbf{Biến định tính:} \texttt{type}, \texttt{rating}, \texttt{language\_name}, \texttt{genres}, \texttt{country}, \texttt{director}, \texttt{cast}.
    \item \textbf{Biến có thứ tự/thời gian:} \texttt{release\_year}, các nhóm điểm chất lượng, nhóm số thể loại.
    \item \textbf{Biến định lượng:} \texttt{vote\_average}, \texttt{vote\_count}, \texttt{popularity}, \texttt{budget}, \texttt{revenue}, ROI dẫn xuất.
\end{itemize}

Đặc điểm quan trọng của dữ liệu là phân phối ngân sách và doanh thu lệch mạnh, nhiều phim có giá trị bằng 0, trong khi điểm đánh giá nằm trong miền hẹp hơn nhưng cần phân nhóm để so sánh chất lượng. Ngoài ra, các trường dạng danh sách như thể loại, đạo diễn và diễn viên buộc phải tách thành cấu trúc ``one-to-many'' để hỗ trợ đếm, xếp hạng và tổng hợp.

\begin{table}[H]
\centering
\renewcommand{\arraystretch}{1.3}
\begin{tabularx}{\textwidth}{|p{3.4cm}|p{2.8cm}|p{3cm}|X|}
\hline
\textbf{Biến} & \textbf{Type} & \textbf{Miền giá trị} & \textbf{Ghi chú cho trực quan hóa} \\
\hline
\texttt{release\_year} & Ordered & 2010--2025 & Phù hợp cho line chart theo thời gian \\
\hline
\texttt{country\_name} & Categorical + Spatial & 133 quốc gia & Phù hợp choropleth để xem mức phủ nội dung \\
\hline
\texttt{language\_name} & Categorical & 68 ngôn ngữ & Treemap phù hợp để thấy cơ cấu tỷ trọng \\
\hline
\texttt{vote\_average} & Quantitative & 0--10 & Có thể bucket hóa để tạo donut/bar/heatmap \\
\hline
\texttt{popularity} & Quantitative & 3.86--3876.01 & Dùng làm thước đo mức độ được quan tâm \\
\hline
\texttt{budget} & Quantitative & 0--460M USD & Phù hợp cho bubble chart, heatmap và quartile \\
\hline
\texttt{revenue} & Quantitative & 0--2.799B USD & Phù hợp so sánh hiệu quả thương mại và ROI \\
\hline
\texttt{genres} & Categorical (list) & Nhiều thể loại/phim & Cần explode để so sánh theo nhóm thể loại \\
\hline
\texttt{director}, \texttt{cast} & Categorical (list) & Nhiều người/phim & Dùng cho xếp hạng và phân cụm sáng tạo \\
\hline
\end{tabularx}
\caption{Bảng phân tích biến số theo What}
\end{table}

\subsection{Why -- Câu hỏi phân tích và nhiệm vụ trực quan}
Các câu hỏi trọng tâm của nhóm được tổ chức thành 4 góc nhìn, bám sát 4 tab của dashboard:
\begin{enumerate}[leftmargin=2em]
    \item \textbf{Toàn cảnh:} Quốc gia nào đóng góp nhiều phim nhất, ngôn ngữ nào chiếm ưu thế, thể loại nào tăng trưởng rõ trong giai đoạn 2010--2025?
    \item \textbf{Chất lượng:} Thư viện phim tập trung ở khoảng điểm nào, thể loại nào vừa có quy mô đủ lớn vừa ổn định về điểm đánh giá?
    \item \textbf{Tài chính:} Mối liên hệ giữa ngân sách, doanh thu và ROI có thực sự tuyến tính hay không; đầu tư đang dịch chuyển vào những thể loại nào?
    \item \textbf{Sáng tạo:} Đạo diễn nào cân bằng tốt giữa doanh thu và điểm số, diễn viên nào kéo doanh thu, và số lượng thể loại của một phim ảnh hưởng thế nào đến mức phổ biến?
\end{enumerate}

Nhóm ánh xạ các câu hỏi trên sang nhiệm vụ trực quan để đảm bảo mỗi biểu đồ đều trả lời ít nhất một task rõ ràng.
\begin{table}[H]
\centering
\renewcommand{\arraystretch}{1.3}
\begin{tabularx}{\textwidth}{|p{4.8cm}|p{3cm}|X|}
\hline
\textbf{Câu hỏi phân tích} & \textbf{Task} & \textbf{Mô tả thao tác người dùng} \\
\hline
Quốc gia nào có độ phủ phim lớn nhất? & Locate / Compare & So sánh trên bản đồ và chọn quốc gia để lọc chéo \\
\hline
Ngôn ngữ nào chiếm tỷ trọng nổi bật? & Part-to-whole & Quan sát treemap để thấy mức tập trung ngôn ngữ \\
\hline
Chất lượng nội dung đang nghiêng về nhóm nào? & Summarize distribution & Quan sát donut và bar theo bucket điểm \\
\hline
Ngân sách có kéo theo doanh thu và ROI? & Find correlation / Compare & Đọc bubble chart, top ROI và line chart theo thời gian \\
\hline
Yếu tố sáng tạo nào nổi bật? & Rank / Explore relationship & Xem scatter đạo diễn, bar diễn viên và biểu đồ kết hợp số thể loại \\
\hline
\end{tabularx}
\caption{Ánh xạ câu hỏi phân tích sang nhiệm vụ trực quan}
\end{table}

\subsection{Phạm vi dashboard}
Dashboard không nhằm giải thích quan hệ nhân quả tuyệt đối, mà tập trung hỗ trợ khám phá mẫu và đặt giả thuyết. Phạm vi của nhóm gồm:
\begin{itemize}[leftmargin=2em]
    \item Quy mô và cấu trúc thư viện phim theo quốc gia, ngôn ngữ và thể loại.
    \item Chất lượng cảm nhận của nội dung thông qua \texttt{vote\_average} và độ ổn định theo thể loại.
    \item Hiệu quả tài chính dựa trên ngân sách, doanh thu và ROI.
    \item Tác động của đạo diễn, diễn viên và độ đa dạng thể loại đến thành công của phim.
    \item Tương tác lọc chéo giữa các biểu đồ để chuyển từ cái nhìn tổng quan sang phân tích chi tiết.
\end{itemize}
